\section{Embeddings in Natural Language Processing (NLP)}

Natural language is an unstructured form of information that is represented through words, sentences, and context rather than explicit numerical structures.
Since machine learning models operate on mathematical representations, NLP systems require a method to transform textual information into a structured format that can be processed computationally. 
Embeddings provide such a representation by mapping words, sentences, or documents into dense vectors within a continuous vector space, where relationships between textual elements can be represented numerically.

The idea of embeddings is based on the distributional hypothesis, which states that words occurring in similar contexts tend to have related meanings~\cite{harris1954}. 
Modern embedding methods use neural networks and large-scale training data to learn vector representations that capture patterns in language usage and contextual relationships.

In contrast to traditional word embeddings, which assign a fixed vector to each word, contextual embedding methods generate representations based on the surrounding text. 
Models such as BERT use transformer architectures to produce context-dependent representations, allowing the same word to receive different representations depending on its usage~\cite{devlin2019bert}.

In this work, embeddings are used to transform unstructured textual information into structured numerical representations that can be compared and processed by mathematical models.